%! Simplifying Rational Expressions — Unit 4, Lesson 4.1 — Group Activity (Answer Key)
\documentclass[10pt]{article}
\usepackage{algebra2-article}
\usepackage{algebra2-key}   % pulls in -boxes; do NOT also load -boxes
\newcommand{\TallMath}[1]{$\displaystyle #1\rule[-1.4em]{0pt}{3.2em}$}
\newcommand{\lgrid}[4]{%
  \draw[step=1, linegray!40, very thin] (#1,#3) grid (#2,#4);
  \draw[->, semithick, charcoal] (#1,0)--(#2,0) node[below right, font=\scriptsize]{$x$};
  \draw[->, semithick, charcoal] (0,#3)--(0,#4) node[left, font=\scriptsize]{$y$};}

\begin{document}

\pageheader{Unit 4, Lesson 4.1}{Group Activity --- Answer Key}
\namepartnerperiod

\vspace{0.08in}

\begin{headlinebox}{forestbg}
\small\textbf{Factor, Restrict, Divide Out.} Every answer today has \emph{two} parts: the simplest
form \textbf{and} the restrictions --- and the restrictions always come from the \textbf{original}
denominator, even when the factor cancels. With your partner, factor first, and be ready to justify
every cancellation by naming the \emph{factor} you divided out.
\end{headlinebox}

\vspace{0.06in}

\begin{tcolorbox}[colback=white, colframe=black!40, title=\textbf{Tier R --- Remediate},
    fonttitle=\bfseries, arc=1.5mm, left=3mm, right=3mm, top=2mm, bottom=2mm, breakable]
Simplify each expression and list the restrictions. The first one is started for you.

\vspace{4pt}
\begin{enumerate}[leftmargin=1.9em, itemsep=9pt]
  \item $\dfrac{10x^4}{15x^2}$: \; the common numerical factor is $5$, and $x^4$ and $x^2$ share
        \ans{$x^2$}.

        Simplest form: \ans{$\frac{2x^2}{3}$} \quad Restriction: $x\neq$ \ans{0}
  \item $\dfrac{x^2-25}{x+5}$ \\[3pt]
        Factored numerator: \ans{$(x-5)(x+5)$} \quad Simplest form: \ans{$x-5$} \quad
        $x\neq$ \ans{$-5$}
  \item $\dfrac{4x+8}{x^2-4}$ \\[3pt]
        Factored: \ans{$\frac{4(x+2)}{(x-2)(x+2)}$} \quad Simplest form: \ans{$\frac{4}{x-2}$} \quad
        $x\neq$ \ans{$2,\ -2$}
  \item $\dfrac{x^2-x-12}{x^2-16}$ \\[3pt]
        Factored: \ans{$\frac{(x-4)(x+3)}{(x-4)(x+4)}$} \quad Simplest form: \ans{$\frac{x+3}{x+4}$} \quad
        $x\neq$ \ans{$4,\ -4$}
  \item In problem 4 the factor $(x-4)$ cancelled. Is $x=4$ still excluded? \ans{Yes} \;
        Why? \ans{the original denominator is $0$ there}
\end{enumerate}
\end{tcolorbox}

\begin{tcolorbox}[colback=white, colframe=black!40, title=\textbf{Tier A --- Approaching Proficiency},
    fonttitle=\bfseries, arc=1.5mm, left=3mm, right=3mm, top=2mm, bottom=2mm, breakable]
\textbf{Part 1 --- The full table.} Fill in every cell. Watch for a GCF, a perfect-square trinomial,
and one pair of \emph{opposite} binomials.
\begin{center}
\renewcommand{\arraystretch}{1.7}
\begin{tabularx}{\linewidth}{@{}l X X X@{}}
\hline
\textbf{Expression} & \textbf{Factored form} & \textbf{Restrictions} & \textbf{Simplest form} \\
\hline
\TallMath{\frac{6x^3y^2}{9xy^5}} & \ans{$\frac{3x y^2\cdot 2x^2}{3xy^2\cdot 3y^3}$} & \ans{$x\neq0,\ y\neq0$} & \ans{$\frac{2x^2}{3y^3}$} \\
\TallMath{\frac{2x^2-18}{x^2+6x+9}} & \ans{$\frac{2(x-3)(x+3)}{(x+3)^2}$} & \ans{$x\neq-3$} & \ans{$\frac{2(x-3)}{x+3}$} \\
\TallMath{\frac{25-x^2}{x^2-3x-10}} & \ans{$\frac{-(x-5)(x+5)}{(x-5)(x+2)}$} & \ans{$x\neq5,\ -2$} & \ans{$-\frac{x+5}{x+2}$} \\
\TallMath{\frac{x^2+2x}{x^2+4x+4}} & \ans{$\frac{x(x+2)}{(x+2)^2}$} & \ans{$x\neq-2$} & \ans{$\frac{x}{x+2}$} \\
\hline
\end{tabularx}
\end{center}

\vspace{4pt}
\textbf{Part 2 --- Error analysis.} A student turns in this work:
\begin{center}
\fbox{\parbox{0.62\linewidth}{\centering \vspace{2pt}
$\dfrac{x^2+9}{x+3} = \dfrac{x^2+9}{x+3} = x+3$ \\[3pt]
{\small ``I cancelled the $9$ with the $3$ and the $x^2$ with the $x$.''}\vspace{2pt}}}
\end{center}
\begin{enumerate}[label=(\alph*), leftmargin=1.9em, itemsep=6pt]
  \item Disprove it with a number. At $x=1$: the original equals \ans{$\frac{10}{4}=2.5$} and the student's
        answer equals \ans{4}.
  \item Name the mistake precisely, using the words \emph{factor} and \emph{term}.
        \ansline{$x^2$ and $9$ are \emph{terms} of the numerator (they are added), not \emph{factors} of it, and}
        \ansline{the same is true of $x$ and $3$ below. Only a common \emph{factor} --- something multiplying the whole numerator and the whole denominator --- may be divided out.}
  \item Can $\dfrac{x^2+9}{x+3}$ be simplified at all? \ans{No} \; Explain what you know about
        $x^2+9$ from Unit 3. \ans{a sum of squares is prime over the reals}
\end{enumerate}
\end{tcolorbox}

\begin{tcolorbox}[colback=white, colframe=black!40, title=\textbf{Tier E --- Extension},
    fonttitle=\bfseries, arc=1.5mm, left=3mm, right=3mm, top=2mm, bottom=2mm, breakable]
\textbf{Part 1 --- Equivalent, but not identical.} Three expressions are offered as ``the same as''
$\dfrac{x-4}{x+1}$.
\begin{center}
\renewcommand{\arraystretch}{1.5}
\begin{tabularx}{\linewidth}{@{}c l X X@{}}
\hline
 & \textbf{Expression} & \textbf{Simplifies to} & \textbf{Extra restriction(s)?} \\
\hline
(i)   & \TallMath{\frac{x^2-16}{x^2+5x+4}} & \ans{$\frac{x-4}{x+1}$} & \ans{yes: $x\neq-4$} \\
(ii)  & \TallMath{\frac{4-x}{x+1}}         & \ans{$-\frac{x-4}{x+1}$} & \ans{none --- but it is \emph{not} equivalent} \\
(iii) & \TallMath{\frac{x^2-4x}{x^2+x}}    & \ans{$\frac{x-4}{x+1}$} & \ans{yes: $x\neq0$} \\
\hline
\end{tabularx}
\end{center}
Two of the three simplify to $\dfrac{x-4}{x+1}$ --- but neither is equal to it at \emph{every} real
input. Name the input each one fails at, and explain why that does not stop them from being
\textbf{equivalent}.
\ansline{(i) fails at $x=-4$ and (iii) fails at $x=0$: each is undefined there while $\frac{x-4}{x+1}$ is not.}
\ansline{Equivalent means equal at every input where \emph{both} are defined, so the mismatch in domain (a hole) does not break equivalence. (ii) is the \emph{opposite} expression --- $-1$ times the target --- so it is not equivalent anywhere.}

\vspace{5pt}
\textbf{Part 2 --- Build one backwards.} Write a rational expression that simplifies to
$\dfrac{x}{x-1}$ but whose restrictions are $x\neq1$ \emph{and} $x\neq-5$.
\begin{center}
\ans{$\frac{x(x+5)}{(x-1)(x+5)}$} \quad Now expand your denominator: \ans{$\frac{x^2+5x}{x^2+4x-5}$}
\end{center}
Explain how you knew where to put the extra factor.
\ansline{the extra restriction must come from the \emph{denominator}, so $(x+5)$ has to appear on the bottom --- and to keep the value unchanged it must appear on top as well, so it cancels}

\vspace{5pt}
\textbf{Part 3 --- Back to the graph.} The graph below belongs to \emph{one} of these expressions.
\begin{center}
\begin{tikzpicture}[scale=0.46]
  \lgrid{-4}{5}{-3}{6}
  \begin{scope}
    \clip (-4,-3) rectangle (5,6);
    \draw[very thick, forest] (-4,-3)--(5,6);
  \end{scope}
  \draw[forest, very thick, fill=white] (3,4) circle (4.2pt);
\end{tikzpicture}
\end{center}
\begin{enumerate}[label=(\alph*), leftmargin=1.9em, itemsep=6pt]
  \item Circle the one it belongs to: \quad
        $\dfrac{x^2-2x-3}{x-3}$ \; \textcolor{keyred}{\textbf{$\leftarrow$ correct}} \qquad
        $\dfrac{x^2-2x-3}{x+1}$ \qquad $\dfrac{x-3}{x^2-2x-3}$
  \item Simplify your choice: \ans{$x+1$}, \; for $x\neq$ \ans{3}. The hole is at
        \ans{$(3,4)$}.
  \item The hole's $y$-value came from the \emph{simplified} expression, not the original. Why can it
        not come from the original?
        \ansline{the original is undefined at $x=3$ --- substituting gives $\frac00$, which is no value at all. The}
        \ansline{simplified form $x+1$ agrees with the original at every \emph{other} input, so it tells us the height the graph was heading for: $4$.}
\end{enumerate}
\end{tcolorbox}

\begin{teachernote}
\textbf{Tier A, Part 1, row 3} is the row that separates groups: $25-x^2$ and $x^2-3x-10$ share the
\emph{opposite} binomials $5-x$ and $x-5$. Accept either $-\frac{x+5}{x+2}$ or $\frac{-(x+5)}{x+2}$ or
$\frac{x+5}{-(x+2)}$, but not a missing sign. Row 2 and row 4 both hide a squared factor --- students
who cancel $(x+3)^2$ entirely will report a restriction-free answer; send them back to the definition
of the factored form.

\vspace{3pt}
\textbf{Tier E, Part 1(ii)} is the deliberate trap: it is the only one that does \emph{not} simplify to
the target, and students who match ``looks similar'' will call it equivalent. Ask them to evaluate both
at $x=0$: $\frac{4}{1}=4$ versus $\frac{-4}{1}=-4$.

\vspace{3pt}
\textbf{Circulate for the two signature errors:} restrictions taken from the \emph{simplified}
denominator (the cancelled restriction vanishes), and cancelling across a $+$ or $-$ sign. When you see
the second, do not correct it verbally --- hand the group a number to test, exactly as in Part 2.
\end{teachernote}

\end{document}
